\begin{table}[H]
  \centering
  \caption{Does an Inc-bind rate hike shift the muni's spending mix? DiD on share$^X$.}
  \label{tab:spending_did}
  \small
  \begin{tabular}{lrrr}
    \toprule
    Outcome (LHS) & Inc-bind $\beta$ & SE & $N$ \\
    \midrule
    Police (exp.) & -0.0143 & (0.0127) & 242,730 \\
    Streets & roads (exp.) & -0.0357 & (0.0298) & 242,730 \\
    Debt service (exp.) & 0.0209 & (0.0251) & 242,730 \\
    EIT (rev. share) & 0.1423$^{***}$ & (0.0275) & 242,730 \\
    Real estate (rev. share) & -0.0709$^{**}$ & (0.0346) & 242,730 \\
    Intergovernmental (rev. share) & -0.0610 & (0.0490) & 242,730 \\
    \bottomrule
  \end{tabular}
  \vspace{0.3em}
  \begin{minipage}{0.95\textwidth}
    \footnotesize \textit{Notes:} Each row is a separate
    continuous-stacked DiD with share$^X$ as the outcome.
    Same FE structure as the headline regression
    (\texttt{stack\_muni + stack\_year}), clustered at
    the muni level, control set is the other five share
    variables. $\beta$ is the Inc-bind dose-response: per
    pp of resident-rate hike at an Inc-bind muni, how much
    the outcome share shifts relative to the NC control
    pool over the same calendar years.
    $^{*}p<0.10$, $^{**}p<0.05$, $^{***}p<0.01$.
  \end{minipage}
\end{table}
